\centerline{\bf \Huge Первая МатДрака }

\problem{\textbf{1}} Чему равно следующее выражение:  $2^{14}-2^{11}-2^{9}+2^{6}-979$


\problem{\textbf{2}} Какое двузначное число увеличивается в 4,5 раза, если его прочитать справа налево?


\problem{\textbf{3}} На каждой из трех полок лежат книги. На одной из них на 22 книги больше, чем на другой. Сколько книг лежит на каждой полке, если всего их 25?


\problem{\textbf{4}} 
Четыре мецената пожертвовали театру 132 тысячи рублей. При этом второй пожертвовал вдвое больше первого, третий — втрое больше второго, четвёртый — вчетверо больше третьего. Сколько пожертвовал четвёртый?


\problem{\textbf{5}} В каждой клетке квадрата 3×3 лежит монета орлом вверх. Какое наименьшее количество монет нужно перевернуть, чтобы в результате не оказалось трех монет, расположенных в одной вертикали, одной горизонтали или одной диагонали и лежащих одинаково (то есть все три орлом вверх или все три решкой вверх)?


\problem{\textbf{6}} 
Ягами выписывает последовательно на доску по возрастанию все числа, в которых число четных цифр равно числу нечетных цифр. Какое число выпишет Ягами 46–м?


\problem{\textbf{7}} 
 АА задумал целое число. Сугуру умножил его не то на 5, не то на 6. Сукуна прибавил к результату Сугуру не то 5, не то 6. Сатору отнял от результата Сукуны не то 5, не то 6. В итоге получилось 73. Какое число задумал АА (перечислите все возможные варианты)?

\problem{\textbf{8}} 
Эмбер пробежала 1 км со средней скоростью 4 м/с. С какой средней скоростью пробежал эту дистанцию Эль Примо, если, стартовав на 25 секунд позже Эмбер, он финишировал на 25 секунд раньше?


\problem{\textbf{9}} Сколько существует трёхзначных чисел, у которых сумма цифр больше произведения цифр?


\problem{\textbf{10}} Сколько на шахматной доске имеется всевозможных прямоугольников, состоящих из четырёх клеток?

\problem{\textbf{11}} У лейтенанта Крипера есть взвод солдатиков. Крипер третью часть своих солдат посадил в засаду, третью часть оставшихся отправил в разведку, а остальные 12 остались рыть окопы. Сколько солдат во взводе?

\newpage
\hspace{6mm}

\problem{\textbf{12}} 
Каждый из 12 человек — рыцарь, всегда говорящий правду, или всегда лгущий лжец. Один из них сказал: «Число лжецов среди нас делится на 1», второй: «Число лжецов среди нас делится на 2», …, 12–ый: «Число лжецов среди нас делится на 12». Сколько среди них может быть рыцарей? Укажите все возможные варианты.


\problem{\textbf{13}} 
Микаса заменила в примере на умножение двузначных чисел цифры буквами: одинаковые – одинаковыми, разные – разными. У неё получилось АБ$\times$ВГ = БББ Каким мог быть исходный пример? Найдите все возможные варианты.


\problem{\textbf{14}} Найдите 2009–й член последовательности: 1, 2, 1, 2, 3, 2, 1, 2, 3, 4, 3, 2, 1, 2, 3, 4, 5, 4, 3, 2, 1…




\problem{\textbf{15}} 
На шахматной доске стоят ладьи так, что каждая из них бьёт $N$ ладей. При каких целых $N$ это возможно? (Ладья бьёт в каждом направлении только ближайшую ладью.) 

